% This is not OFFICIAL TEMPLATE OF COEPTech UNI
% compile with XeLaTeX
\documentclass[dvipsnames,mathserif,aspectratio=169,10pt]{beamer}
%you can change aspect ratio
\setbeamertemplate{footline}[frame number]
\setbeamercolor{footline}{fg=black}
\setbeamerfont{footline}{series=\bfseries}
\usepackage{tikz}
\usepackage{xcolor}
\usepackage{amsmath}


\usetheme{Frankfurt}%1


% for RTL liste
\makeatletter
\newcommand{\RTListe}{\raggedleft\rightskip\leftm}
\newcommand{\leftm}{\@totalleftmargin}
\makeatother


% RTL frame title
\setbeamertemplate{frametitle}
{\vspace*{-1mm}
  \nointerlineskip
    \begin{beamercolorbox}[sep=0.3cm,ht=2.2em,wd=\paperwidth]{frametitle}
        \vbox{}\vskip-2ex%
        \strut\hskip1ex\insertframetitle\strut
        \vskip-0.8ex%
    \end{beamercolorbox}
}
% align subsection in toc
\makeatletter
\setbeamertemplate{subsection in toc}
{\leavevmode\rightskip=5ex%
  \llap{\raise0.1ex\beamer@usesphere{subsection number projected}{bigsphere}\kern1ex}%
  \inserttocsubsection\par%
}
\makeatother
%LOGO
\logo{%
  \makebox[0.98\paperwidth,]{
    \includegraphics[width=0.80cm,keepaspectratio]{Logos/COEP-Tech-Pune-Logo-3013335431.png} \hspace{5cm}
    \includegraphics[scale=0.20]{Logos/co.png}
    \hfill
    \includegraphics[]{Logos/Logo.png}
    
  }
}
%\logo{\includegraphics[height=1cm]{COEP-Tech-Pune-Logo-3013335431.png}}
% RTL triangle for itemize
\setbeamertemplate{itemize item}{\scriptsize\raise1.25pt\hbox{\donotcoloroutermaths$\blacktriangleleft$}} 

%\setbeamertemplate{itemize item}{\rule{4pt}{4pt}}

\defbeamertemplate{enumerate item}{square2}
{\LR{
    %
    \hbox{%
    \usebeamerfont*{item projected}%
    \usebeamercolor[bg]{item projected}%
    \vrule width2.25ex height1.85ex depth.4ex%
    \hskip-2.25ex%
    \hbox to2.25ex{%
      \hfil%
      {\color{fg}\insertenumlabel}%
      \hfil}%
  }%
}}

\setbeamertemplate{enumerate item}[square2]
\setbeamertemplate{navigation symbols}{}
% title
\title[About] %optional
{Comparative Analysis of Graph Traversal and
Shortest Path Algorithms}

% name, mis & affiliation
\author[Sumit] % optional short name 
{Sumit Desai}
\institute[COEP] 
{
    \text{Department of Computer Science \& Engineering} \\
    \text{COEP Technological University} \\
    Pune, India \\
    MIS: 612415049
}

\begin{document}

\begin{frame}
\maketitle
\transfade[duration=0.2]
\end{frame}


%
 \footnotesize \tableofcontents   
%
\section{Abstract} 
\begin{frame}{Abstract}

    \begin{itemize}
        \item \textbf{Context:} Models pairwise relationships in complex systems (e.g., social networks, transportation).
        \item \textbf{Mathematical Definition:} Systems are modeled as a graph $G$ defined in Eq.~\eqref{eq:graphdef}.
    \end{itemize}

    \begin{equation}
        G = (V, E) \label{eq:graphdef}
    \end{equation}
    {\small Where $V$ represents vertices and $E$ represents edges.}

    \vspace{0.5em}

    \textbf{Algorithms Analyzed:}
    \begin{enumerate}
        \item Breadth-First Search (BFS)
        \item Depth-First Search (DFS)
        \item Dijkstra's Algorithm
    \end{enumerate}
\end{frame}


\section{Introduction}
\begin{frame}{Introduction}
\begin{columns}

\column{0.5\textwidth}
\textbf{Background and Motivation}
\begin{itemize}
    \item Graphs are widely used to model real-world systems such as networks, maps, and social platforms.
    \item Efficient traversal of graphs is essential for solving connectivity and routing problems.
    \item Early approaches focused on simple traversal without cost considerations.
    \item Traditional methods often ignored edge weights or scalability issues.
    \item Growing data sizes demand more efficient and optimal algorithms.
    \item This motivates the study of classical and shortest-path algorithms.
\end{itemize}

\column{0.5\textwidth}
\textbf{Problem Overview}

Graph traversal and shortest path problems form the foundation of many computer science applications. While simple traversal techniques are sufficient for unweighted graphs, they fail to account for real-world constraints such as distance, cost, or time. This creates a need to analyze and compare different algorithms based on their behavior, efficiency, and suitability for specific graph types. Understanding these differences helps in selecting the right algorithm for practical applications.

\end{columns}
\transfade[duration=0.2]
\end{frame}


\section{Methodology}
\begin{frame}{Methodology: Breadth-First Search (BFS)}
    \begin{columns}
        \column{0.5\textwidth}
        \begin{itemize}
            \item Earlier traversal methods focused only on visiting nodes, not path optimality.
            \item BFS improves this by exploring graphs level by level.
            \item It guarantees the shortest path only when all edges have equal cost.
            \item BFS fails in weighted graphs because it ignores edge weights.
            \item This limitation demands more cost-aware approaches → Dijkstra.
        \end{itemize}
    \column{0.5\textwidth}
        \begin{figure}
            \includegraphics[height=0.50\linewidth]{bfs_demo.jpeg}
            \caption*{Level-wise expansion from source node}
        \end{figure}
    \end{columns}
\transfade[duration=0.2]
\end{frame}


\begin{frame}{Methodology: Depth-First Search (DFS)}
    \begin{columns}
        \column{0.5\textwidth}
        \begin{itemize}
            \item DFS explores deeply before considering alternative paths.
            \item Earlier approaches using DFS were efficient for structure analysis.
            \item DFS does not guarantee shortest or optimal paths.
            \item It may traverse long paths even when shorter paths exist.
            \item Hence, DFS is unsuitable for routing or cost-based problems.
        \end{itemize}
    \column{0.5\textwidth}
        \begin{figure}
            \includegraphics[height=0.50\linewidth]{dfs_demo.jpeg}
            \caption*{Deep traversal along a single branch}
        \end{figure}
    \end{columns}
\transfade[duration=0.2]
\end{frame}

\begin{frame}{Methodology: Dijkstra's Algorithm}
    \begin{columns}
        \column{0.5\textwidth}
            \begin{itemize}
                \item BFS and DFS fail to handle weighted graphs correctly.
                \item Dijkstra introduces a greedy, cost-based selection strategy.
                \item It always expands the node with minimum cumulative cost.
                \item Uses a priority queue to ensure optimal path selection.
                \item Overcomes limitations of earlier traversal-based approaches.
            \end{itemize}
            
        \column{0.5\textwidth}
            \begin{figure}
                \centering
                \includegraphics[height=0.50\linewidth]{dijkstras_demo.jpeg}
                \caption*{Shortest path based on cumulative edge weight}
            \end{figure}
    
    \end{columns}
\transfade[duration=0.2]
\end{frame}

\section{Results}
\begin{frame}{Results}
    \begin{itemize}
        \item The comparative study shows that no single graph algorithm is optimal for all scenarios.
        \item BFS performs best for shortest path problems in unweighted graphs with minimal overhead.
        \item DFS is effective for structural analysis but unsuitable for optimal path discovery.
        \item Dijkstra's algorithm produces correct shortest paths in weighted graphs at a higher computational cost.
        \item Future work includes extending this study to algorithms handling negative weights and large-scale graphs.
    \end{itemize}
    \transfade[duration=0.2]
\end{frame}

\begin{frame}{Algorithm Comparison Summary}
\begin{table}
\centering
\begin{tabular}{|l|l|l|l|}
\hline
\textbf{Metric} & \textbf{BFS} & \textbf{DFS} & \textbf{Dijkstra} \\
\hline
Graph Type & Unweighted & Any & Weighted \\
\hline
Shortest Path & Yes & No & Yes \\
\hline
Time Complexity & $O(V+E)$ & $O(V+E)$ & $O(E\log V)$ \\
\hline
Auxiliary Space & $O(V)$ & $O(V)$ & $O(V)$ \\
\hline
Primary Use & Traversal & Structure & Routing \\
\hline
\end{tabular}
\end{table}
\transfade[duration=0.2]
\end{frame}



\section{Conclusion}
\begin{frame}{Conclusion}
    \begin{itemize}
        \item A detailed literature review of classical graph algorithms was conducted.
        \item BFS, DFS, and Dijkstra's algorithm were analyzed theoretically and visually.
        \item Algorithm behavior was studied using traversal diagrams and complexity analysis.
        \item Comparative results highlight strengths and limitations of each approach.
        \item Future work includes implementing simulations and analyzing advanced shortest path algorithms.
    \end{itemize}
    \transfade[duration=0.2]
\end{frame}



\section{References}
\begin{frame}{References}
    \begin{itemize}
        \item [1] T. H. Cormen et al., \textit{Introduction to Algorithms}, MIT Press.
        \item [2] E. W. Dijkstra, ``A note on two problems in connexion with graphs,'' 1959.
        \item [3]  R. Sedgewick and K. Wayne, \textit{Algorithms}, Addison-Wesley.
        \item [4] TutorialsPoint, ``Graph traversal algorithms,'' Online resource.
    \end{itemize}
\transfade[duration=0.2]
\end{frame}

\end{document}